problem:  
Let \((X,d,m)\) be a noncollapsed \(\mathrm{RCD}(0,N)\) space with \(N<\infty\).  
For \(m\)-a.e. \(x \in X\), the tangent cone is of the canonical form  
\[
T_xX \simeq \mathbb{R}^k \times C(Y_x),
\]
where \(Y_x\) is an \(\mathrm{RCD}(N-k-1,N-k)\) space with \(\operatorname{diam}(Y_x)\le \pi\).  
Assume that for some fixed compact \(\mathrm{RCD}(N-k-1,N-k)\) space \(Y\) we have  
\[
T_xX \simeq \mathbb{R}^k \times C(Y)
\quad\text{for } m\text{-a.e. }x\in X.
\]
Assume moreover that the **density function**
\[
\vartheta(x) = \lim_{r\to 0} \frac{m(B_r(x))}{\omega_N r^N}
\]
exists and is **continuous** on a dense open subset \(U\subset X\).  

**Problem:**  
Does it follow that \(U\) is locally biLipschitz (with quantitative control depending only on \(K,N\)) to a product \(\mathbb{R}^k\times Z\),  
where \(Z\) is an \(\mathrm{RCD}(0,N-k)\) space that admits tangents \(C(Y)\) at \(m\)-a.e. point?  
Equivalently, does having constant cross‑section \(Y\) and continuous volume density enforce the **existence of local product charts** modeled on \(\mathbb{R}^k\times C(Y)\)?  

Why is it a "good" problem:  
This formulation refines the microlocal‑to‑local rigidity question to one that is both **precisely defined** and **genuinely open** in current RCD theory.  

* **Precisely stated objective:** It asks for local biLipschitz product coordinates—explicit geometric regularity stronger than rectifiability but weaker than global splitting. The meaning of "locally product‑like" is now fully specified.  
* **Grounded in existing structure theory:** It builds directly on Mondino–Naber (rectifiability of regular strata) and De Philippis–Gigli (uniqueness of tangents a.e.), yet it extends these by demanding *alignment and coherence* of the cone factor across points, which is not covered by current theorems.  
* **Subtle and nontrivial:** Continuity of density ensures measure‑theoretic regularity, but it is unknown whether this suffices to integrate infinitesimal cones into local product neighborhoods. This confronts the core unsolved problem of gluing identical tangents into local geometry, parallel to open questions in Cheeger–Colding theory for Ricci limit spaces.  
* **Conceptually deep:** The question tests whether constancy of infinitesimal structure (cone model) and absence of density jumps enforce local geometric rigidity—an analogue of “volume cone ⇒ metric cone” but in the infinitesimal‑to‑local regime.  
* **Potential impact:** A positive answer would substantiate a synthetic version of local splitting and could lead toward classification results for noncollapsed RCD spaces with simple singular structures; a negative example would reveal new topological or geometric obstructions to aligning tangent geometry.  

Its statement is simple, yet resolution would demand new analytic–geometric techniques uniting local measure regularity, blow‑up stability, and curvature‑dimension rigidity—the hallmarks of a truly “good” research problem.